\bigskip\Needspace{10\baselineskip}
\pdfbookmark[1]{Further Report. Ceres Ferdinandea. March 1802}{ceres180203}
\begin{center}
{\small FURTHER REPORT ON THE NEW MAJOR PLANET.\par}
CERES FERDINANDEA.\par
\medskip\hrule\medskip
Monatliche Correspondenz. Vol.\ V. pp.\ 263—282. March 1802.\par
\medskip\hrule
\end{center}\medskip
—\quad —\quad —\quad “I cannot help mentioning here, writes this distinguished scholar [\textsc{Gauss}], what a benefit it has been to astronomy on this occasion to have a journal like the M.\ C.\ How lukewarmly and indifferently would \textsc{Piazzi}'s discovery have been received if you had not gathered all the reports about it in your journal, circulated them as quickly as possible, aroused general interest, weighed the arguments for and against, and made it as probable as could be that this object was a planet. Probably only a few astronomers would have taken the trouble to look for it again, since even the teacher and master of all present-day astronomers had so strongly doubted the new planet only a short time before?” —\quad —\quad —

—\quad —\quad In the December issue of last year, p.\ 646 [p.\ 203 of this volume], we gave our readers all the orbital elements that Dr.\ \textsc{Gauss} had calculated in succession and improved up to the \textit{fourth time} on the basis of \textsc{Piazzi}'s observations. Dr.\ \textsc{Gauss} later made the following small changes to this \textit{fourth} set of elements:
\begin{quote}
$\AscendingNode\ 81^0\ 0'\ 13''39$, aphelion $326^0\ 27'\ 41''93$, epoch $77^0\ 36'\ 38''35$, inclination $10^0\ 37'\ 1''98$.
\end{quote}
The rest remained as it was. —\quad —

—\quad —\quad —\quad In the preceding February issue, p.\ 178, we already mentioned that before \textit{Ceres} was found, Dr.\ \textsc{Gauss} had sent us an orbit of this planet improved for the \textit{fifth time}. So we wish for now to acquaint our astronomical readers with this orbit as well, for the historical record. Dr.\ \textsc{Gauss} subjected \textsc{Piazzi}'s observations to a repeated, careful calculation of the elliptical orbit, taking the most precise account of nutation, aberration and parallax. This yielded the following fifth set of elements:
\SourceLeaf{211}%
\begin{center}
(V)\par\smallskip
{\small\setlength{\tabcolsep}{2pt}
\begin{tabular}{l r}
Aphelion $\ .\ .\ .\ .\ .\ .\ .\ .\ .$&$324^0\ 37'\ 11''$\\
$\AscendingNode\ .\ .\ .\ .\ .\ .\ .\ .\ .\ .\ .\ .$&$80\ 59\ 12$\\
Inclination $\ .\ .\ .\ .\ .\ .\ .\ .\ .\ .$&$10\ 37\ 9''55$\\
Logarithm of the semimajor axis&$0.4446804$\\
Eccentricity $\ .\ .\ .\ .\ .\ .\ .\ .\ .$&$0.0879111$\\
mean daily heliocentric mean motion&$763''950$\\
Epoch 1800 31 December $\ .\ .\ .\ .$&$78^0\ 5'\ 16''6$
\end{tabular}}
\end{center}
It would be superfluous to show here how well these elements agree with \textsc{Piazzi}'s observations, which Dr.\ \textsc{Gauss} had calculated with the greatest care and represented to within a couple of seconds.

When I told Dr.\ \textsc{Gauss} the good news that the long-awaited \textit{Ceres} had been found and sent him my first three observations of it, the first thing he did on receiving them was to calculate them at once from his V\textsuperscript{th} elements given above. For the first observation, on 7 December, he found an error in right ascension of $+24'\ 8''$; for the second, on 11 January, one of $+30'\ 53''$; for the third, on 16 January, one of $+31'\ 53''$. By a rough calculation with his IV\textsuperscript{th} (improved) elements, the deviations in the same order are $+14\frac12$ minutes $+19\frac34$ minutes $+20\frac12$ minutes. Dr.\ \textsc{Gauss} takes it to be a coincidence that these IV elements come somewhat closer to the truth than the V\textsuperscript{th}; perhaps it is also partly due to the effect of planetary perturbations on \textsc{Piazzi}'s observations, especially on the latitudes. As far as Dr.\ \textsc{Gauss} could tell from the first two observations by \textsc{Olbers}, which he had received in the meantime, the latitudes agreed with both sets of elements to within a couple of minutes, and the difference in declinations was mainly just a consequence of the differences in longitude. Before receiving further observations from me, however, Dr.\ \textsc{Gauss} could not resist the temptation to make a provisional correction to his elements. He used my two right ascensions of 7 December and 16 January and thus found the following VI\textsuperscript{th} elements after a calculation carried out with all possible care. Such care was all the more necessary since the slightest error would have made the whole work useless.
\begin{center}
{\small\setlength{\tabcolsep}{2pt}
\begin{tabular}{l l}
Epoch 1801 Palermo meridian $\ .\ .$&$77^0\ 24'\ 55''9$\\
—\quad1802\quad—\quad— $\ .\ .$&$155\ 33\ 35.1$\\
Mean daily tropical motion&$770''7376$\\
Tropical period $\ .\ .\ .\ .\ .\ .$&1681 days 12 hr. 9 min.\\
Logarithm of the semimajor axis&$0.4421189$
\end{tabular}
\par\smallskip
$\left.\begin{array}{lr}
\text{Aphelion}\ .\ .\ .\ .\ .\ .\ .\ .\ .&326^0\ 14'\ 45''\\
\AscendingNode\ .\ .\ .\ .\ .\ .\ .\ .\ .\ .\ .\ .&80\ 58\ 55
\end{array}\right\}\quad\begin{gathered}
\text{31 December 1800,}\\
\text{assumed fixed with respect to the stars}
\end{gathered}$
\par\smallskip
\begin{tabular}{l r}
Eccentricity $\ .\ .\ .\ .\ .\ .\ .\ .\ .$&$0.08086253$\\
greatest equation of the centre $\ .\ .\ .$&$9^0\ 16'\ 23''$\\
Inclination $\ .\ .\ .\ .\ .\ .\ .\ .\ .\ .\ .$&$10\ 37\ 51$
\end{tabular}}
\end{center}
It was in the nature of the method used in this calculation that these elements represent \textsc{Piazzi}'s observations almost as precisely as the V\textsuperscript{th} elements. But it was not to be expected that this first improvement would bring the errors in the two new observations from half a degree down to zero at once. Still, when Dr.\ \textsc{Gauss} had calculated from these new elements the two right ascensions on which they were based, he found the agreement closer than he had hoped. The error in the first observation was only $+3''0$, and in the last it was $+9''9$. Dr.\ \textsc{Gauss} had also calculated from these elements three new observations of mine that I had sent him in the meantime, and he found the following agreement in right ascension.

\newcommand{\CeresSymbol}{\text{\mdseries\AstroSymbols ⚳}}
\SourceLeaf{212}%
\begin{center}
{\small\setlength{\tabcolsep}{5pt}
\begin{tabular}{r r|r|r|r}
\multicolumn{2}{c|}{Seeberg}&\shortstack{Calculated\\R.A. $\CeresSymbol$}&Difference&\shortstack{Calculated\\declination}\\\hline
1801 December&7&$178^0\ 33'\ 33''6$&$+3''0$&$11^0\ 47'\ 33''$ N\\
1802 January&11&$186\ 46\ 9.3$&$+19.3$&$11\ 15\ 41$\\
—&; 16&$187\ 28\ 3.1$&$+9.9$&$11\ 26\ 40$\\
—&22&$188\ 6\ 45.9$&$+20.1$&$11\ 45\ 18$\\
—&25&$188\ 21\ 6.5$&$+27.3$&$11\ 56\ 49$
\end{tabular}}
\end{center}
Since the only latitudes used to determine these elements were \textsc{Piazzi}'s, all of which lie very close to the node, we cannot yet expect great precision in the inclination. Still, after a rough calculation using \textsc{Olbers}'s declinations, Dr.\ \textsc{Gauss} believes the inclination is likely to be in error by scarcely more than a minute. This is also further evidence of the excellent quality of \textsc{Piazzi}'s declinations.

These VI\textsuperscript{th} elements will still need considerable improvement later. The farther apart the observations are, the more certain the improvements will be than if Dr.\ \textsc{Gauss} were already to fit his elements as closely as possible to the observations now available. Yet Dr.\ \textsc{Gauss} is so skilled and tireless a calculator that finding new elements seems to cost him no effort at all. Before I knew it, and before I had sent him my observation of 26 January, he surprised me with a VII\textsuperscript{th} orbit of \textit{Ceres}. For this he used, in place of my observation of the 16th, the one of 25 January and so obtained the following (VII) elements.
\begin{center}
(VII)
\par\smallskip
{\small
$\text{Epoch 1801}\ \left\{\begin{array}{lr}
\text{Seeberg meridian}\ .\ .&77^0\ 27'\ 36''5\\
\text{Palermo meridian}\ .\ .&77\ 27\ 30.9
\end{array}\right.$
\par\smallskip
\begin{tabular}{l r}
mean daily tropical motion $\ .$&$769''7924$\\
Logarithm of the semimajor axis $\ .$&$0.4424742$\\
Eccentricity $\ .\ .\ .\ .\ .\ .\ .\ .\ .\ .$&$0.0814064$
\end{tabular}
\par\smallskip
$\left.\begin{array}{l}
\text{Aphelion}\\
\AscendingNode
\end{array}\right\}\ \text{1801 fixed}\ \begin{array}{lr}
.\ .\ .\ .\ .&325^0\ 57'\ 15''\\
.\ .\ .\ .\ .&80\ 58\ 40
\end{array}$
\par\smallskip
\begin{tabular}{l r}
greatest equation of the centre $\ .\ .\ .$&$9\ 20\ 8.0$\\
Inclination $\ .\ .\ .\ .\ .\ .\ .\ .\ .\ .\ .$&$10\ 37\ 56.6$\\
{}[Handwritten] Epoch 1802 Palermo $\ .\ .$&$155\ 30\ 25.2$
\end{tabular}}
\end{center}
Here again only \textsc{Piazzi}'s latitudes were used. These elements also represent all the Palermo observations as precisely as possible, and now agree more closely with my observations as well, as the following comparison shows:
\begin{center}
{\small\setlength{\tabcolsep}{4pt}
\begin{tabular}{r r|r|r|r|r}
\multicolumn{2}{c|}{Seeberg}&\shortstack{Calculated\\R.A.}&Difference&\shortstack{Calculated\\declination N.}&Difference\\\hline
1801 December&7&$178^0\ 33'\ 29''2$&$-1''4$&&\\
1802 January&11&$186\ 45\ 47.6$&$-2.3$&&\\
&16&$187\ 27\ 38.8$&$-14.4$&&\\
&22&$188\ 6\ 18.2$&$-7.6$&&\\
&25&$188\ 20\ 37.2$&$-2.0$&$11^0\ 56'\ 58''4$&$+35''4$\\
&26&$188\ 24\ 37.0$&$-12.5$&$12\quad1\quad8.4$&\\
&28&$188\ 31\ 25.7$&$-12.1$&$12\quad9\ 55.6$&$+14.3$\\
&29&$188\ 34\ 14.1$&$-4.0$&$12\ 14\ 34.1$&\\
&30&$188\ 36\ 38.4$&$-5.5$&$12\ 19\ 19.8$&$+19.1$\\
&31&$188\ 38\ 38.3$&$-7.1$&$12\ 24\ 15.3$&\\
February&3&$188\ 42\quad9.5$&$-3.5$&$12\ 39\ 56.6$&$+17.6$\\
&4&$188\ 42\ 30.7$&$-5.6$&$12\ 45\ 25.7$&\\
&5&$188\ 42\ 26.7$&$-3.4$&$12\ 51\quad2.6$&$+34.7$\\
&9&$188\ 38\quad1.6$&$-2.3$&$13\ 14\ 43.0$&$+22.0$
\end{tabular}}
\end{center}
\expandafter\def\csname HistoricalNoteText18\endcsname{To leave nothing undone here either that might help us observe this planet for as long as possible and attain the greatest possible accuracy, we will now, while there is still time, try to determine very precisely in advance all the stars that will lie in that part of the apparent orbit of \textit{Ceres} where the planet can no longer be observed in the meridian and can be observed at its western quadrature only with equatorial instruments or parallactic telescopes. These observations also deserve greater attention because they will be the most favorable for determining the \textit{radius vector}.}
\SourceLeaf{213}%
It would be easy for Dr.\ \textsc{Gauss} to make the small difference above between the calculation and all my observed right ascensions to date disappear entirely. But we can now hope that improving these elements will give us a greater degree of accuracy almost every week, and in any case they are accurate enough to find \textit{Ceres} for a long time. So Dr.\ \textsc{Gauss} himself believes it best for now to continue comparing the observations with \textit{one and the same set of elements} for a while longer and, in particular, to wait for a somewhat larger number of accurate declinations. Even the small discrepancy shown by the table above, together with the insufficient number of reliable declinations, already shows that it is too soon for a fur-further correction of the VII elements. Moreover, the orbit of \textit{Ceres} is no perfect ellipse because of the perturbations it suffers from the action of the other planets. These perturbations must first be calculated before the \textit{true} elliptical orbit can be obtained. —\quad —\quad —\quad Nevertheless, as things now stand, with \textit{Ceres} having already traversed almost a quarter of its orbit, Dr.\ \textsc{Gauss} flatters himself that the difference between these VII elements and the true elements of the orbit is likely to be small. With all the observations we may still hope for this year before \textit{Ceres} becomes invisible\HistoricalNote{18}, he hopes to bring these elements to such a degree of accuracy that, when one calculates back 44 years with them to \textsc{Mayer}'s observations, one should be unable to miss by as much as a degree in longitude and should miss by only a few minutes in latitude. If \textit{Ceres} should fail to turn up in any older star catalogue as \textit{Uranus} did, we will of course learn its mean motion more accurately only by degrees, and for a long time there will still be something to add and improve each year. —\quad —\quad —

To help our astronomical readers find \textit{Ceres} in the future as well, we give them here another short \textit{ephemeris} of its course for the next two months, which Dr.\ \textsc{Gauss} calculated from his latest and VII elements. In April these positions will hopefully differ from the sky by no more than a minute.
\begin{center}
Positions of \textit{Ceres} in March and April 1802 at midnight, mean Seeberg time.
\par\smallskip
{\small\setlength{\tabcolsep}{5pt}
\begin{tabular}{r r|r|r|r}
\multicolumn{2}{c|}{1802}&\shortstack{R.A. $\CeresSymbol$\\in degrees}&Decl.&R.A. in time\\\hline
March&1&$186^0\ 41'$&$15^0\ 30'$ N&$12^{\mathrm u}\ 26^{\mathrm m}\ 45^{\mathrm s}$\\
&4&$186\ 11$&$15\ 50$&$12\ 24\ 45$\\
&7&$185\ 39$&$16\ 10$&$12\ 22\ 36$\\
&10&$185\quad5$&$16\ 29$&$12\ 20\ 18$\\
&13&$184\ 28$&$16\ 47$&$12\ 17\ 53$\\
&16&$183\ 51$&$17\quad4$&$12\ 15\ 24$\\
&19&$183\ 13$&$17\ 19$&$12\ 12\ 50$\\
&22&$182\ 34$&$17\ 33$&$12\ 10\ 15$\\
&25&$181\ 55$&$17\ 44$&$12\quad7\ 40$
\end{tabular}
\par\smallskip
\begin{tabular}{r r|r|r|r}
\multicolumn{2}{c|}{1802}&\shortstack{R.A. $\CeresSymbol$\\in degrees}&Decl.&R.A. in time\\\hline
March&28&$181^0\ 17'$&$17^0\ 54'$ N&$12^{\mathrm u}\quad5^{\mathrm m}\quad7^{\mathrm s}$\\
&31&$180\ 39$&$18\quad1$&$12\quad2\ 37$\\
April&3&$180\quad3$&$18\quad6$&$12\quad0\ 12$\\
&6&$179\ 29$&$18\quad9$&$11\ 57\ 54$\\
&9&$178\ 56$&$18\ 10$&$11\ 55\ 45$\\
&12&$178\ 26$&$18\quad9$&$11\ 53\ 46$\\
&15&$177\ 59$&$18\quad5$&$11\ 51\ 57$\\
&18&$177\ 35$&$17\ 59\frac12$&$11\ 50\ 20$
\end{tabular}}
\end{center}
Dr.\ \textsc{Gauss} adds to these ephemerides the following important remarks derived from them:

“The planet's opposition falls on the afternoon of 17 March. That will be the best time to observe it. Shortly before opposition it has reached its least distance from the Earth
\SourceLeaf{214}%
$=1.6025$. At about the same time (a little earlier) it attains its greatest northern geocentric latitude $=17^0\ 9'$, and a little later its fastest retrograde motion, about 13 min. of right ascension per day. The northern declination continues to increase until the beginning of April; about 9 April the motion will turn south again.” —\quad —\quad —
\begin{flushright}[\textsc{von Zach}.]\end{flushright}

